\documentclass[11pt, a4paper]{article}
\usepackage{graphicx}
\usepackage[T2A]{fontenc}
\usepackage[utf8]{inputenc}
\usepackage[russian]{babel}
\usepackage{amsmath, amssymb}
\usepackage[usenames]{color}
\usepackage{colortbl}
\usepackage{fancyhdr}
\usepackage{enumitem}
\usepackage{titlesec}
\usepackage{multicol}
\usepackage{tabularx}
\usepackage{caption}
\usepackage{tikz}
\usetikzlibrary{positioning, arrows.meta}
\usepackage[backend=biber, style=numeric, sorting=none]{biblatex}
\DefineBibliographyExtras{russian}{\protected\def\bibrangedash{\textendash\penalty\hyphenpenalty}}
\AtBeginBibliography{\sloppy}
\renewcommand*{\bibfont}{\small}
\addbibresource{source.bib}
\usepackage{csquotes}
\usepackage[hidelinks]{hyperref}
\usepackage{xurl}
\usepackage{placeins}

\usepackage{geometry}
\geometry{
  a4paper,
  top=30mm,
  right=5mm,
  bottom=20mm,
  left=5mm,
  headheight=61pt
}
\setlength{\columnsep}{8mm}
\raggedcolumns  % не растягивать колонки по высоте
\emergencystretch=2em
\renewcommand{\topfraction}{0.95}
\renewcommand{\bottomfraction}{0.6}
\renewcommand{\textfraction}{0.05}
\renewcommand{\floatpagefraction}{0.85}
\setcounter{topnumber}{3}\setcounter{totalnumber}{4}

\newcommand{\mydate}{Октябрь, 2026}
\newcommand{\myauthor}{В. Алексеев, С. Зинина, И. Коновалов, В. Горбатенко, А. Соколова}

\fancyhf{}
\fancyhead[L]{\large \textit{Научная студия 2026. Вычислительная онкология. Нормализация данных и кодирование признаков.}}
\fancyfoot[C]{\thepage}
\pagestyle{fancy}

\fancypagestyle{firstpage}{
  \fancyhf{}
  \fancyhead[L]{\hspace{-15pt}\includegraphics[width=0.3\linewidth]{fig/cu_logo.png}}
  \fancyhead[R]{\mydate\\\textit{\myauthor}\\}
  \fancyfoot[C]{\thepage}
}

\newcommand{\code}[1]{\texttt{\small #1}}

\begin{document}
\thispagestyle{firstpage}

\begin{multicols}{2}
\section*{Материалы и методы}

\subsection*{Разделение данных}

Аналитическая таблица (140 пациентов, 12 с рецидивом) разделена на обучающую и тестовую выборки стратифицированно по исходу с фиксированным начальным значением генератора случайных чисел, поэтому при повторном запуске получаются те же выборки \cite{Pedregosa2011}. Стратификация нужна потому, что при 12 событиях случайное разделение могло бы оставить в тестовой выборке одно событие или ни одного. В тестовую выборку отведена четверть пациентов: 35, из них 3 с рецидивом. В обучающую вошли 105 пациентов, из них 9 с рецидивом. Большая доля теста сократила бы и без того малое число событий для обучения \cite{Riley2020}, а меньшая оставила бы для итоговой оценки одно-два события.\vspace{6pt}\\
Тестовая выборка будет использована один раз, для итоговой оценки выбранной модели. Подбор моделей и их параметров будет выполняться перекрёстной проверкой внутри обучающей выборки. Дополнительной внешней проверкой станут предсказания для 158 пациентов тестовой части соревнования на Kaggle \cite{KaggleOPC}. Все величины для подготовки признаков (значения для заполнения пропусков, перечень логарифмируемых признаков, сдвиг и масштаб, категории для кодирования) вычислены только по обучающей выборке без использования исхода и без пересчёта применены к тестовой.

\subsection*{Подготовка признаков}

В обучающей выборке пропусков нет у 79 из 105 пациентов. Степень дифференцировки пропущена у 13 пациентов, пачко-годы у 9, минимальная плотность у 5, а ещё 11 признаков плотности и текстуры~--- у одного пациента, у которого они удалены из-за металла в маске опухоли. Исключение пациентов с пропусками сократило бы выборку на четверть, поэтому пропуски заполнены. Все пропуски пачко-лет приходятся на курящих и бывших курильщиков, и они заполнены медианой среди курильщиков обучающей выборки (26 пачко-лет): общая медиана была бы занижена некурящими с нулевым значением. Степень дифференцировки заполнена медианой порядкового кода (III). Отсутствие оценки само может нести сведения о биопсийном материале, поэтому добавлен признак-индикатор пропуска. Радиомические признаки заполнены медианами обучающей выборки, так как медиана устойчива к асимметрии распределений.\vspace{6pt}\\
Выбросы в каждом количественном признаке обучающей выборки определялись по правилу полутора межквартильных размахов. Они найдены у 17 из 20 признаков, больше всего у дозы за фракцию (20), максимальной плотности (11), <<занятости>> текстуры (busyness, 10) и объёма опухоли (9). Выбросы не удалялись и не усекались. Значения, невозможные для тканей, исправлены на этапе первичного анализа, а оставшиеся отражают реальные случаи: альтернативные схемы фракционирования и крупные опухоли. Их удаление сократило бы и без того малое число пациентов с рецидивом. Влияние выбросов ослаблено логарифмированием восьми признаков с правой асимметрией больше 1 и неотрицательными значениями: $\ln x$ для строго положительных и $\ln(1+x)$ для пачко-лет, где есть нули. Логарифм $\ln(1+x)$ для признаков со значениями порядка сотых почти не меняет распределение, поэтому для них взят $\ln x$. После логарифмирования выбросов не осталось в объёме опухоли, пачко-годах и <<занятости>> текстуры, а в максимальной плотности их стало 1 вместо 11. Признаки с левой асимметрией, например 10-й перцентиль плотности, логарифмированию не подлежат и сохранены как есть.\vspace{6pt}\\
Количественные признаки нормализованы вычитанием среднего и делением на стандартное отклонение обучающей выборки (z-нормализация). Далее планируются модели с регуляризацией, штраф которых зависит от масштаба признаков, поэтому признаки приведены к сопоставимому масштабу. После логарифмирования распределения ближе к симметричным, и среднее со стандартным отклонением описывают их более чем удовлетворительно. Так же нормализованы порядковые коды.\vspace{6pt}\\
Бинарные признаки (пол, индукционная и одновременная химиотерапия) закодированы значениями 0 и 1. Порядковые признаки закодированы рангами, сохраняющими порядок: категория T~--- от 1 до 4, категория N (0, 1, 2a, 2b, 2c, 3)~--- от 0 до 5, стадия~--- от 1 до 4, степень дифференцировки (I, I--II, II, II--III, III, IV)~--- значениями 1; 1,5; 2; 2,5; 3; 4. Прямое кодирование отдельными столбцами потеряло бы порядок и добавило бы признаков при малом числе событий. Номинальные признаки (раса, сторона и подобласть опухоли, статус курения) закодированы бинарными столбцами для каждой категории, кроме опорной (One-Hot-Encoding), самой частой в обучающей выборке: европеоидная раса, левая сторона, корень языка, бывший курильщик. Опорная категория исключена, чтобы столбцы не были линейно зависимы. Категории, встречающиеся в обучающей выборке реже 5 раз, объединены, так как отдельный столбец для 1--3 пациентов даёт неустойчивую оценку. Так, негроидная и латиноамериканская расы объединены в другую расу, а мягкое нёбо и стенка глотки отнесены к другой подобласти. Категории, отсутствующие в обучающей выборке, при кодировании тестовой выборки также относятся к другой.\vspace{6pt}\\
В итоге подготовленные выборки содержат 35 признаков: 20 количественных, 4 порядковых, индикатор пропуска степени дифференцировки, 3 бинарных и 7 бинарных столбцов номинальных признаков. Пропусков не осталось, все признаки числовые.

\end{multicols}

\FloatBarrier
\begin{multicols}{2}
\subsection*{Ссылки}
\printbibliography[heading=none]
\end{multicols}

\end{document}
